\section*{Resumen}
\addcontentsline{toc}{section}{Resumen}
Se estudia la clasificación binaria de ventas unitarias mensuales de Wide World Importers (WWI), a partir de una extracción de su Data Mart de Ventas. La respuesta indica si las unidades observadas de un producto superan la mediana de sus tres meses previos. Se auditaron 228\,265 líneas de factura y se construyeron observaciones producto-mes con información disponible antes del horizonte objetivo. El perceptrón simple y la regresión logística se implementaron mediante NumPy sobre las mismas 20 características transformadas, con particiones temporales de 4\,599 observaciones de entrenamiento, 2\,628 de validación y 1\,111 de prueba.

La selección utilizó únicamente validación; las predicciones finales se conservaron y su consistencia se comprobó sin reentrenar. En prueba, la regresión logística alcanzó balanced accuracy de 62,42\%, frente a 54,10\% del perceptrón, 50,00\% del baseline mayoritario y 51,31\% de la regla histórica. La logística registró 370 falsos positivos y 49 falsos negativos, lo que limita su utilidad para alertas automáticas. El perceptrón mantuvo actualizaciones hasta la última época; la logística mostró descenso estable de entropía cruzada, sin demostración de convergencia completa al agotar 200 épocas. PCA se empleó solo para visualización y conservó 49,62\% de la varianza en dos componentes.

La ampliación financiera estudia una respuesta empresarial distinta: frecuencia de cobros registrados por cuenta cliente superior a la mediana de tres meses anteriores. A partir de eventos operacionales tipo 3 se prepararon 16 características y 4\,612/2\,943/1\,284 observaciones de entrenamiento/validación/prueba. En la población principal de prueba, el perceptrón alcanzó BA de \TFBAPTest\% y logística \TFBALTest\%, con mayor sensibilidad y más falsas alertas en el primero. El histórico contractual requiere cuatro meses y se compara por separado sobre 1\,276 filas. El PCA financiero, solo visual, conserva \TFPCARetenida\% de varianza en dos componentes. Ninguno de los historiales demuestra convergencia completa.

Las conclusiones son descriptivas de cinco meses de prueba por aplicación, con observaciones dependientes y cobertura académica, no certificada operacionalmente. Ventas conserva identidad provisional por nombre; Transacciones usa CustomerID operacional y supuestos explícitos de conservación del libro de eventos. Este informe unifica ambas aplicaciones sin comparar sus porcentajes como si correspondieran a la misma etiqueta o población. Pedidos y Compras permanecen como investigaciones de viabilidad con limitaciones documentadas; Inventario/Movimientos y la revisión de los procesos bloqueados forman parte del trabajo académico pendiente.

\noindent\textbf{Palabras clave:} Inteligencia de Negocios; Data Warehouse; clasificación temporal; perceptrón; regresión logística; PCA; trazabilidad.
